\begin{textsample}{POS Dim 1 – 1990 – Score 4.00 – 1990/t000287.txt}  \label{ex:f1_pos_075}
I ' m going to go right into \textbf{the} \textbf{slides}.
And all I ' m going to try and prove to you with these \textbf{slides} is that I do just very straight stuff.
And my ideas are — in my head, anyway — they ' re very logical and relate to what ' s going on and problem solving for clients.
I either convince clients at \textbf{the} end that I solve their problems, or I really do solve their problems, because usually they seem to like it.
Let me go right into \textbf{the} \textbf{slides}.
Can you turn off \textbf{the} light?
Down.
I like to be in \textbf{the} dark.
I don ' t want you to see what I ' m doing up here.
Anyway, I did this house in Santa Monica, and it got a lot of notoriety.
In fact, it appeared in a porno comic book, which is \textbf{the} \textbf{slide} on \textbf{the} right.
This is in Venice.
I just show it because I want you to know I ' m concerned about context.
On \textbf{the} left-hand side, I had \textbf{the} context of those little houses, and I tried to build a building that fit into that context.
When people take pictures of these buildings out of that context they look really weird, and my premise is that they make a lot more sense when they ' re photographed or seen in that space.
And then, once I deal with \textbf{the} context, I then try to make a place that ' s comfortable and private and fairly serene, as I hope you ' ll find that \textbf{slide} on \textbf{the} right.
And then I did a law school for Loyola in downtown L.
A.
I was concerned about making a place for \textbf{the} study of law.
And we continue to work with this client. \textbf{The} building on \textbf{the} right at \textbf{the} top is now under construction. \textbf{The} garage on \textbf{the} right — \textbf{the} gray structure — will be torn down, finally, and several small classrooms will be placed along this avenue that we ' ve created, this campus.
And it all related to \textbf{the} clients and \textbf{the} students from \textbf{the} very first meeting saying they felt denied a place.
They wanted a sense of place.
And so \textbf{the} whole idea here was to create that kind of space in downtown, in a neighborhood that was difficult to fit into.
And it was my theory, or my point of view, that one didn ' t upstage \textbf{the} neighborhood — one made accommodations.
I tried to be inclusive, to include \textbf{the} buildings in \textbf{the} neighborhood, whether they were buildings I liked or not.
In \textbf{the} ' 60s I started working with paper furniture and made a bunch of stuff that was very successful in Bloomingdale ' s.
We even made flooring, walls and everything, out of cardboard.
And \textbf{the} success of it threw me for a loop.
I couldn ' t deal with \textbf{the} success of furniture — I wasn ' t secure enough as an architect — and so I closed it all up and made furniture that nobody would like.
So, nobody would like this.
And it was in this, preliminary to these pieces of furniture, that Ricky and I worked on furniture by \textbf{the} slice.
And after we failed, I just kept failing. \textbf{The} piece on \textbf{the} left — and that ultimately led to \textbf{the} piece on \textbf{the} right — happened when \textbf{the} kid that was working on this took one of those long strings of stuff and folded it up to put it in \textbf{the} wastebasket.
And I put a piece of tape around it, as you see there, and realized you could sit on it, and it had a lot of resilience and strength and so on.
So, it was an accidental discovery.
I got into \textbf{fish}.
I mean, \textbf{the} story I tell is that I got mad at postmodernism — at po-mo — and said that \textbf{fish} were 500 million years earlier than man, and if you ' re going to go back, we might as well go back to \textbf{the} beginning.
And so I started making these funny things.
And they started to have a life of their own and got bigger — as \textbf{the} one glass at \textbf{the} Walker.
And then, I sliced off \textbf{the} head and \textbf{the} tail and everything and tried to translate what I was learning about \textbf{the} form of \textbf{the} \textbf{fish} and \textbf{the} movement.
And a lot of my architectural ideas that came from it — accidental, again — it was an intuitive kind of thing, and I just kept going with it, and made this proposal for a building, which was only a proposal.
I did this building in Japan.
I was taken out to dinner after \textbf{the} contract for this little restaurant was signed.
And I love sake and Kobe and all that stuff.
And after I got — I was really drunk — I was asked to do some sketches on napkins.
And I made some sketches on napkins — little boxes and Morandi-like things that I used to do.
And \textbf{the} client said,"Why no \textbf{fish}?"And so I made a drawing with a \textbf{fish}, and I left Japan.
Three weeks later, I received a complete set of drawings saying we ' d won \textbf{the} competition.
Now, it ' s hard to do.
It ' s hard to translate a \textbf{fish} form, because they ' re so beautiful — perfect — into a building or \textbf{object} like this.
And Oldenburg, who I work with a little once in a while, told me I couldn ' t do it, and so that made it even more exciting.
But he was right — I couldn ' t do \textbf{the} tail.
I started to get \textbf{the} head OK, but \textbf{the} tail I couldn ' t do.
It was pretty hard. \textbf{The} thing on \textbf{the} right is a snake form, a ziggurat.
And I put them together, and you walk between them.
It was a dialog with \textbf{the} context again.
Now, if you saw a picture of this as it was published in Architectural Record — they didn ' t show \textbf{the} context, so you would think,"God, what a pushy guy this is."But a friend of mine spent four hours wandering around here looking for this restaurant.
Couldn ' t find it.
So...
As for craft and technology and all those things that you ' ve all been talking about, I was thrown for a complete loop.
This was built in six months. \textbf{The} way we sent drawings to Japan : we used \textbf{the} magic \textbf{computer} in Michigan that does carved models, and we used to make foam models, which that thing scanned.
We made \textbf{the} drawings of \textbf{the} \textbf{fish} and \textbf{the} scales.
And when I got there, everything was perfect — except \textbf{the} tail.
So, I decided to cut off \textbf{the} head and \textbf{the} tail.
And I made \textbf{the} \textbf{object} on \textbf{the} \textbf{left} for my show at \textbf{the} Walker.
And it ' s one of \textbf{the} nicest pieces I ' ve ever made, I think.
And then Jay Chiat, a friend and client, asked me to do his headquarters building in L.
A.
For reasons we don ' t want to talk about, it got delayed.
Toxic waste, I guess, is \textbf{the} key clue to that one.
And so we built a temporary building — I ' m getting good at temporary — and we put a conference room in that ' s a \textbf{fish}.
And, finally, Jay dragged me to my hometown, Toronto, Canada.
And there is a story — it ' s a real story — about my grandmother buying a carp on Thursday, bringing it home, putting it in \textbf{the} bathtub when I was a kid.
I played with it in \textbf{the} evening.
When I went to sleep, \textbf{the} next day it wasn ' t there.
And \textbf{the} next night, we had gefilte \textbf{fish}.
And so I set up this interior for Jay ' s offices and I made a pedestal for a sculpture.
And he didn ' t buy a sculpture, so I made one.
I went around Toronto and found a bathtub like my grandmother ' s, and I put \textbf{the} \textbf{fish} in.
It was a joke.
I play with funny people like [ Claes ] Oldenburg.
We ' ve been friends for a long time.
And we ' ve started to work on things.
A few years ago, we did a performance piece in Venice, Italy, called"Il Corso del Coltello"— \textbf{the} Swiss Army knife.
And most of \textbf{the} imagery is — Claes ', but those two little boys are my sons, and they were Claes ' assistants in \textbf{the} play.
He was \textbf{the} Swiss Army knife.
He was a souvenir salesman who always wanted to be a painter, and I was Frankie P.
Toronto.
P for Palladio.
Dressed up like \textbf{the} AT\&T building by Claes — with a \textbf{fish} hat. \textbf{The} highlight of \textbf{the} performance was at \textbf{the} end.
This beautiful \textbf{object}, \textbf{the} Swiss Army knife, which I get credit for participating in.
And I can tell you — it ' s totally an Oldenburg.
I had nothing to do with it. \textbf{The} only thing I did was, I made it possible for them to turn those blades so you could sail this thing in \textbf{the} canal, because I love sailing.
We made it into a sailing craft.
I ' ve been known to mess with things like chain link fencing.
I do it because it ' s a curious thing in \textbf{the} culture, when things are made in such great quantities, absorbed in such great quantities, and there ' s so much denial about them.
People hate it.
And I ' m fascinated with that, which, like \textbf{the} paper furniture — it ' s one of those materials.
And I ' m always drawn to that.
And so I did a lot of dirty things with chain link, which nobody will forgive me for.
But Claes made homage to it in \textbf{the} Loyola Law School.
And that chain link is really expensive.
It ' s in perspective and everything.
And then we did a camp together for children with cancer.
And you can see, we started making a building together.
Of course, \textbf{the} milk can is his.
But we were trying to collide our ideas, to put \textbf{objects} next to each other.
Like a Morandi — like \textbf{the} little bottles — composing them like a still life.
And it seemed to work as a way to put he and I together.
Then Jay Chiat asked me to do this building on this funny lot in Venice, and I started with this three-piece thing, and you entered in \textbf{the} middle.
And Jay asked me what I was going to do with \textbf{the} piece in \textbf{the} middle.
And he pushed that.
And one day I had a — oh, well, \textbf{the} other way.
I had \textbf{the} binoculars from Claes, and I put them there, and I could never get rid of them after that.
Oldenburg made \textbf{the} binoculars incredible when he sent me \textbf{the} first model of \textbf{the} real proposal.
It made my building look sick.
And it was this interaction between that kind of, up-the-ante stuff that became pretty interesting.
It led to \textbf{the} building on \textbf{the} \textbf{left}.
And I still think \textbf{the} Time magazine picture will be of \textbf{the} binoculars, you know, leaving out \textbf{the} — what \textbf{the} hell.
I use a lot of metal in my work, and I have a hard time connecting with \textbf{the} craft. \textbf{The} whole thing about my house, \textbf{the} whole use of rough carpentry and everything, was \textbf{the} frustration with \textbf{the} crafts available.
I said,"If I can ' t get \textbf{the} craft that I want, I ' ll use \textbf{the} craft I can get."There were plenty of models for that, in Rauschenberg and Jasper Johns, and many artists who were making beautiful art and sculpture with junk materials.
I went into \textbf{the} metal because it was a way of building a building that was a sculpture.
And it was all of one material, and \textbf{the} metal could go on \textbf{the} roof as well as \textbf{the} walls. \textbf{The} metalworkers, for \textbf{the} most part, do ducts behind \textbf{the} ceilings and stuff.
I was given an opportunity to design an exhibit for \textbf{the} metalworkers ' unions of America and Canada in Washington, and I did it on \textbf{the} condition that they become my partners in \textbf{the} future and help me with all future metal buildings, etc. etc.
And it ' s working very well to have these people, these craftsmen, interested in it.
I just tell \textbf{the} stories.
It ' s a way of connecting, at least, with some of those people that are so important to \textbf{the} realization of architecture. \textbf{The} metal continued into a building — Herman Miller, in Sacramento.
And it ' s just a complex of factory buildings.
And Herman Miller has this philosophy of having a place — a people place.
I mean, it ' s kind of a trite thing to say, but it is real that they wanted to have a central place where \textbf{the} cafeteria would be, where \textbf{the} people would come and where \textbf{the} people working would interact.
So it ' s out in \textbf{the} middle of nowhere, and you approach it.
It ' s copper and galvanize.
I used \textbf{the} galvanize and copper in a very light gauge, so it would buckle.
I spent a lot of time undoing Richard Meier ' s aesthetic.
Everybody ' s trying to get \textbf{the} panels perfect, and I always try to get them sloppy and fuzzy.
And they end up looking like stone.
This is \textbf{the} central area.
There ' s a ramp.
And that little dome in there is a building by Stanley Tigerman.
Stanley was instrumental in my getting this job.
And when I was awarded \textbf{the} contract I, at \textbf{the} very beginning, asked \textbf{the} client if they would let Stanley do a cameo piece with me.
Because these were ideas that we were talking about, building things next to each other, making — it ' s all about [ a ] metaphor for a city, maybe.
And so Stanley did \textbf{the} little dome thing.
And we did it over \textbf{the} phone and by fax.
He would send me a fax and show me something.
He ' d made a building with a dome and he had a little tower.
I told him,"No, no, that ' s too ongepotchket.
I don ' t want \textbf{the} tower."So he came back with a simpler building, but he put some funny details on it, and he moved it closer to my building.
And so I decided to put him in a depression.
I put him in a hole and made a kind of a hole that he sits in.
And so then he put two bridges — this all happened on \textbf{the} fax, going back and forth over a couple of weeks ' period.
And he put these two bridges with pink guardrails on it.
And so then I put this big billboard behind it.
And I call it,"David and Goliath."And that ' s my cafeteria.
In Boston, we had that old building on \textbf{the} \textbf{left}.
It was a very prominent building off \textbf{the} freeway, and we added a floor and cleaned it up and fixed it up and used \textbf{the} kind of — I thought — \textbf{the} language of \textbf{the} neighborhood, which had these cornices, projecting cornices.
Mine got a little exuberant, but I used lead copper, which is a beautiful material, and it turns green in 100 years.
Instead of, like, copper in 10 or 15.
We redid \textbf{the} side of \textbf{the} building and re-proportioned \textbf{the} windows so it sort of fit into \textbf{the} space.
And it surprised both Boston and myself that we got it approved, because they have very strict kind of design guideline, and they wouldn ' t normally think I would fit them. \textbf{The} detailing was very careful with \textbf{the} lead copper and making panels and fitting it tightly into \textbf{the} fabric of \textbf{the} existing building.
In Barcelona, on Las Ramblas for some film festival, I did \textbf{the} Hollywood sign going and coming, made a building out of it, and they built it.
I flew in one night and took this picture.
But they made it a third smaller than my model without telling me.
And then more metal and some chain link in Santa Monica — a little shopping center.
And this is a laser laboratory at \textbf{the} University of Iowa, in which \textbf{the} \textbf{fish} comes back as an abstraction in \textbf{the} back.
It ' s \textbf{the} support labs, which, by some coincidence, required no windows.
And \textbf{the} shape fit perfectly.
I just joined \textbf{the} points.
In \textbf{the} curved part there ' s all \textbf{the} mechanical equipment.
That solid wall behind it is a pipe chase — a pipe canyon — and so it was an opportunity that I seized, because I didn ' t have to have any protruding ducts or vents or things in this form.
It gave me an opportunity to make a sculpture out of it.
This is a small house somewhere.
They ' ve been building it so long I don ' t remember where it is.
It ' s in \textbf{the} West Valley.
And we started with \textbf{the} stream and built \textbf{the} house along \textbf{the} stream — dammed it up to make a lake.
These are \textbf{the} models. \textbf{The} reality, with \textbf{the} lake — \textbf{the} workmanship is pretty bad.
And it reminded me why I play defensively in things like my house.
When you have to do something really cheaply, it ' s hard to get perfect corners and stuff.
That big metal thing is a passage, and in it is — you go downstairs into \textbf{the} living room and then down into \textbf{the} bedroom, which is on \textbf{the} right.
It ' s kind of like a whole built town.
I was asked to do a hospital for schizophrenic adolescents at Yale.
I thought it was fitting for me to be doing that.
This is a house next to a Philip Johnson house in Minnesota. \textbf{The} owners had a dilemma — they asked Philip to do it.
He was too busy.
He didn ' t recommend me, by \textbf{the} way.
We ended up having to make it a sculpture, because \textbf{the} dilemma was, how do you build a building that doesn ' t look like \textbf{the} language?
Is it going to look like this beautiful estate is sub-divided?
Etc. etc.
You ' ve got \textbf{the} idea.
And so we finally ended up making it.
These people are art collectors.
And we finally made it so it appears very sculptural from \textbf{the} main house and all \textbf{the} windows are on \textbf{the} other side.
And \textbf{the} building is very sculptural as you walk around it.
It ' s made of metal and \textbf{the} brown stuff is Fin-Ply — it ' s that formed lumber from Finland.
We used it at Loyola on \textbf{the} chapel, and it didn ' t work.
I keep trying to make it work.
In this case we learned how to detail it.
In Cleveland, there ' s Burnham Mall, on \textbf{the} \textbf{left}.
It ' s never been finished.
Going out to \textbf{the} lake, you can see all those new buildings we built.
And we had \textbf{the} opportunity to build a building on this site.
There ' s a railroad track.
This is \textbf{the} city hall over here somewhere, and \textbf{the} courthouse.
And \textbf{the} centerline of \textbf{the} mall goes out.
Burnham had designed a railroad station that was never built, and so we followed.
Sohio is on \textbf{the} axis here, and we followed \textbf{the} axis, and they ' re two kind of goalposts.
And this is our building, which is a corporate headquarters for an insurance company.
We collaborated with Oldenburg and put \textbf{the} newspaper on top, folded. \textbf{The} health club is fastened to \textbf{the} garage with a C-clamp, for Cleveland.
You drive down.
So it ' s about a 10-story C-clamp.
And all this stuff at \textbf{the} bottom is a museum, and an idea for a very fancy automobile entry.
This owner has a pet peeve about bad automobile entries.
And this would be a hotel.
So, \textbf{the} centerline of this thing — we ' d preserve it, and it would start to work with \textbf{the} scale of \textbf{the} new buildings by Pelli and Kohn Pederson Fox, etc., that are underway.
It ' s hard to do high-rise.
I feel much more comfortable down here.
This is a piece of property in Brentwood.
And a long time ago, about ' 82 or something, after my house — I designed a house for myself that would be a village of several pavilions around a courtyard — and \textbf{the} owner of this lot worked for me and built that actual model on \textbf{the} \textbf{left}.
And she came back, I guess wealthier or something — something happened — and asked me to design a house for her on this site.
And following that basic idea of \textbf{the} village, we changed it as we got into it.
I locked \textbf{the} house into \textbf{the} site by cutting \textbf{the} back end — here you see on \textbf{the} photographs of \textbf{the} site — slicing into it and putting all \textbf{the} bathrooms and dressing rooms like a retaining wall, creating a lower level zone for \textbf{the} master bedroom, which I designed like a kind of a barge, looking like a boat.
And that ' s it, built. \textbf{The} dome was a request from \textbf{the} client.
She wanted a dome somewhere in \textbf{the} house.
She didn ' t care where.
When you sleep in this bedroom, I hope — I mean, I haven ' t slept in it yet.
I ' ve offered to marry her so I could sleep there, but she said I didn ' t have to do that.
But when you ' re in that room, you feel like you ' re on a kind of barge on some kind of lake.
And it ' s very private. \textbf{The} landscape is being built around to create a private garden.
And then up above there ' s a garden on this side of \textbf{the} living room, and one on \textbf{the} other side.
These aren ' t focused very well.
I don ' t know how to do it from here.
Focus \textbf{the} one on \textbf{the} right.
It ' s up there.
Left — it ' s my right.
Anyway, you enter into a garden with a beautiful grove of trees.
That ' s \textbf{the} living room.
Servants ' quarters.
A guest bedroom, which has this dome with marble on it.
And then you enter into \textbf{the} living room and then so on.
This is \textbf{the} bedroom.
You come down from this level along \textbf{the} stairway, and you enter \textbf{the} bedroom here, going into \textbf{the} lake.
And \textbf{the} bed is back in this space, with windows looking out onto \textbf{the} lake.
These Stonehenge things were designed to give foreground and to create a greater depth in this shallow lot. \textbf{The} material is lead copper, like in \textbf{the} building in Boston.
And so it was an intent to make this small piece of land — it ' s 100 by 250 — into a kind of an estate by separating these areas and making \textbf{the} living room and dining room into this pavilion with a high space in it.
And this happened by accident that I got this right on axis with \textbf{the} dining room table.
It looks like I got a Baldessari painting for free.
But \textbf{the} idea is, \textbf{the} windows are all placed so you see pieces of \textbf{the} house outside.
Eventually this will be screened — these trees will come up — and it will be very private.
And you feel like you ' re in your own kind of village.
This is for Michael Eisner — Disney.
We ' re doing some work for him.
And this is in Anaheim, California, and it ' s a freeway building.
You go under this bridge at about 65 miles an hour, and there ' s another bridge here.
And you ' re through this room in a split second, and \textbf{the} building will sort of reflect that.
On \textbf{the} backside, it ' s much more humane — entrance, dining hall, etc.
And then this thing here — I ' m hoping as you drive by you ' ll hear \textbf{the} picket fence effect of \textbf{the} sound hitting it.
Kind of a fun thing to do.
I ' m doing a building in Switzerland, Basel, which is an office building for a furniture company.
And we struggled with \textbf{the} image.
These are \textbf{the} early studies, but they have to sell furniture to normal people, so if I did \textbf{the} building and it was too fancy, then people might say,"Well, \textbf{the} furniture looks OK in his thing, but no, it ain ' t going to look good in my normal building."So we ' ve made a kind of pragmatic slab in \textbf{the} second phase here, and we ' ve taken \textbf{the} conference facilities and made a villa out of them so that \textbf{the} communal space is very sculptural and separate.
And you ' re looking at it from \textbf{the} offices and you create a kind of interaction between these pieces.
This is in Paris, along \textbf{the} Seine.
Palais des Sports, \textbf{the} Gare de Lyon over here. \textbf{The} Minister of Finance — \textbf{the} guy that moved from \textbf{the} Louvre — goes in here.
There ' s a new library across \textbf{the} river.
And back in here, in this already treed park, we ' re doing a very dense building called \textbf{the} American Center, which has a theater, apartments, dance school, an art museum, restaurants and all kinds of — it ' s a very dense program — bookstores, etc.
In a very tight, small — this is \textbf{the} ground level.
And \textbf{the} French have this extraordinary way of screwing things up by taking a beautiful site and cutting \textbf{the} corner off.
They call it \textbf{the} plan coupe.
And I struggled with that thing — how to get around \textbf{the} corner.
These are \textbf{the} models for it.
I showed you \textbf{the} other model, \textbf{the} one — this is \textbf{the} way I organized myself so I could make \textbf{the} drawing — so I understood \textbf{the} problem.
I was trying to get around this plan coupe — how do you do it?
Apartments, etc.
And these are \textbf{the} kind of study models we did.
And \textbf{the} one on \textbf{the} \textbf{left} is pretty awful.
You can see why I was ready to commit suicide when this one was built.
But out of it came finally this resolution, where \textbf{the} elevator piece worked frontally to this, parallel to this street, and also parallel to here.
And then this kind of twist, with this balcony and \textbf{the} skirt, kind of like a ballerina lifting her skirt to let you into \textbf{the} foyer. \textbf{The} restaurants here — \textbf{the} apartments and \textbf{the} theater, etc.
So it would all be built in stone, in French limestone, except for this metal piece.
And it faces into a park.
And \textbf{the} idea was to make this express \textbf{the} energy of this.
On \textbf{the} side facing \textbf{the} street it ' s much more normal, except I slipped a few mansards down, so that coming on \textbf{the} point, these housing units made a gesture to \textbf{the} corner.
And this will be some kind of high-tech billboard.
If any of you guys have any ideas for it, please contact me.
I don ' t know what to do.
Jay Chiat is a glutton for punishment, and he hired me to do a house for him in \textbf{the} Hamptons.
And it ' s got a \textbf{fish}.
And I keep thinking,"This is going to be \textbf{the} last \textbf{fish}."It ' s like a drug addict.
I say,"I ' m not going to do it anymore — I don ' t want to do it anymore — I ' m not going to do it."And then I do it.
There it is.
But it ' s \textbf{the} living room.
And this piece here is — I don ' t know what it is.
I just added it so that we ' d have enough money in \textbf{the} budget so we could take something out.
This is Euro Disney, and I ' ve worked with all of \textbf{the} guys that presented to you earlier.
We ' ve had a lot of fun working together.
I think I ' m from Mars for them, and they are for me, but somehow we all manage to work together, and I think, productively.
So far.
This is a shopping thing.
You come into \textbf{the} Magic Kingdom and \textbf{the} hotel that Tony Baxter ' s group is doing out here.
And then this is a kind of a shopping mall, with a rodeo and restaurants.
And another restaurant.
What I did — because of \textbf{the} Paris skies being quite dull, I made a light grid that ' s perpendicular to \textbf{the} train station, to \textbf{the} route of \textbf{the} train.
It looks like it ' s kind of been there, and then crashed all these simpler forms into it. \textbf{The} light grid will have a light, be lit up at night and give a kind of light ceiling.
In Switzerland — Germany, actually — on \textbf{the} Rhine across from Basel, we did a furniture factory and a furniture museum.
And I tried to — there ' s a Nick Grimshaw building over here, there ' s an Oldenburg sculpture over here — I tried to make a relationship urbanistically.
And I don ' t gave good \textbf{slides} to show — it ' s just been completed — but this piece here is this building, and these pieces here and here.
And as you pass by it ' s always part — you see it as all of these pieces accrue and become part of an overall neighborhood.
It ' s plaster and just zinc.
And you wonder, if this is a museum, what it ' s going to be like inside?
If it ' s going to be so busy and crazy that you wouldn ' t show anything, and just wait.
I ' m so cunning and clever — I made it quiet and wonderful.
But on \textbf{the} outside it does scream out at you a bit.
It ' s actually basically three square rooms with a couple of skylights and stuff.
And from \textbf{the} building in \textbf{the} back, you see it as an iceberg floating by in \textbf{the} hills.
I know I ' m over time.
See, that skylight goes down and becomes that one.
So it ' s pretty quiet inside.
This is \textbf{the} Disney Hall — \textbf{the} concert hall.
It ' s a complicated project.
It has a chamber hall.
It ' s related to an existing Chandler Pavilion that was built with a lot of love and tears and caring.
And it ' s not a great building, but I approached it optimistically, that we would make a compositional relationship between us that would strengthen both of us.
And \textbf{the} plan of this — it ' s a concert hall.
This is \textbf{the} foyer, which is kind of a garden structure.
There ' s commercial at \textbf{the} ground floor.
These are offices, which, really, in \textbf{the} competition, we didn ' t have to design.
But finally, there ' s a hotel there.
These were \textbf{the} kind of relationships made to \textbf{the} Chandler, composing these elevations together and relating them to \textbf{the} buildings that existed — to MOCA, etc. \textbf{The} acoustician in \textbf{the} competition gave us criteria, which led to this compartmentalized scheme, which we found out after \textbf{the} competition would not work at all.
But everybody liked these forms and liked \textbf{the} space, and so that ' s one of \textbf{the} problems of a competition.
You have to then try and get that back in some way.
And we studied many models.
This was our original model.
These were \textbf{the} three buildings that were \textbf{the} ideal — \textbf{the} Concertgebouw, Boston and Berlin.
Everybody liked \textbf{the} surround.
Actually, this is \textbf{the} smallest hall in size, and it has more seats than any of these because it has double balconies.
Our client doesn ' t want balconies, so — and when we met our new acoustician, he told us this was \textbf{the} right shape or this was \textbf{the} right shape.
And we tried many shapes, trying to get \textbf{the} energy of \textbf{the} original design within an acoustical, acceptable format.
We finally settled on a shape that was \textbf{the} proportion of \textbf{the} Concertgebouw with \textbf{the} sloping outside walls, which \textbf{the} acoustician said were crucial to this and later decided they weren ' t, but now we have them.
And our idea is to make \textbf{the} seating carriage very sculptural and out of wood and like a big boat sitting in this plaster room.
That ' s \textbf{the} idea.
And \textbf{the} corners would have skylights and these columns would be structural.
And \textbf{the} nice thing about introducing columns is they give you a kind of sense of proscenium from wherever you sit, and create intimacy.
Now, this is not a final design — these are just on \textbf{the} way to being — and so I wouldn ' t take it literally, except \textbf{the} feeling of \textbf{the} space.
We studied \textbf{the} acoustics with laser stuff, and they bounce them off this and see where it all works.
But you get \textbf{the} sense of \textbf{the} hall in section.
Most halls come straight down into a proscenium.
In this case we ' re opening it back up and getting skylights in \textbf{the} four corners.
And so it will be quite a different shape. \textbf{The} original building, because it was frog-like, fit nicely on \textbf{the} site and cranked itself well.
When you get into a box, it ' s harder to do it — and here we are, struggling with how to put \textbf{the} hotel in.
And this is a teapot I designed for Alessi.
I just stuck it on there.
But this is how I do work.
I do take pieces and bits and look at it and struggle with it and cut it away.
And of course it ' s not going to look like that, but it is \textbf{the} crazy way I tend to work.
And then finally, in L.
A.
I was asked to do a sculpture at \textbf{the} foot of Interstate Bank Tower, \textbf{the} highest building in L.
A.
Larry Halprin is doing \textbf{the} stairs.
And I was asked to do a \textbf{fish}, and so I did a snake.
It ' s a public space, and I made it kind of a garden structure, and you can go in it.
It ' s a kiva, and Larry ' s putting some water in there, and it works much better than a \textbf{fish}.
In Barcelona I was asked to do a \textbf{fish}, and we ' re working on that, at \textbf{the} foot of a Ritz-Carlton Tower being done by Skidmore, Owings and Merrill.
And \textbf{the} Ritz-Carlton Tower is being designed with exposed steel, non-fire proof, much like those old gas tanks.
And so we took \textbf{the} language of this exposed steel and used it, perverted it, into \textbf{the} form of \textbf{the} \textbf{fish}, and created a kind of a 19th- century contraption that looks like, that will sit — this is \textbf{the} beach and \textbf{the} harbor out in front, and this is really a shopping center with department stores.
And we split these bridges.
Originally, this was all solid with a hole in it.
We cut them loose and made several bridges and created a kind of a foreground for this hotel.
We showed this to \textbf{the} hotel people \textbf{the} other day, and they were terrified and said that nobody would come to \textbf{the} Ritz-Carlton anymore, because of this \textbf{fish}.
And finally, I just threw these in — Lou Danziger.
I didn ' t expect Lou Danziger to be here, but this is a building I did for him in 1964, I think.
A little studio — and it ' s sadly for sale.
Time goes on.
And this is my son working with me on a small fast-food thing.
He designed \textbf{the} robot as \textbf{the} cashier, and \textbf{the} head moves, and I did \textbf{the} rest of it.
And \textbf{the} food wasn ' t as good as \textbf{the} stuff, and so it failed.
It should have been \textbf{the} other way around — \textbf{the} food should have been good first.
It didn ' t work.
Thank you very much.

% matched lemmas: computer, fish, left, object, slide, the
\end{textsample}
